\originalchapter{宇宙常数、因果过去与超新星习题}
本章完成8.286 PS8五题。数值积分、共同因果条件求根及60行历史超新星图均已实际执行。代码与输入来源校验记录随源码保存；它们分别对应此题指定的背景模型、资料和误差假设，不能用8.902的SNLS拟合替代。

\src{8.286 PS8.1，亮度与红移，(a)--(b)。}
\begin{exercise}
今天给$H_0,\Omega_{m0},\Omega_{r0},\Omega_{v0}$，$\Omega_{k0}=1-\Omega_{m0}-\Omega_{r0}-\Omega_{v0}$。已知闭宇宙的$J$与$z$关系，(a)完整推导开宇宙版本；(b)推平坦版本，若取曲率趋零极限，要处理分子、分母同时趋零。
\end{exercise}
\begin{solution}
用$a_0=1$的尺度因子，$1+z=a^{-1}$，定义
\[
 E(z)=\frac{H(z)}{H_0}
 =\sqrt{\Omega_{m0}(1+z)^3+\Omega_{r0}(1+z)^4+
 \Omega_{v0}+\Omega_{k0}(1+z)^2},\quad
 I(z)=\int_0^z\frac{\dd z'}{E(z')}.
\]
径向零测地线给共动光路$\chi=(c/H_0)I$。对开宇宙$\Omega_{k0}>0$，今天曲率半径$R_{K0}=c/(H_0\sqrt{\Omega_{k0}})$，其横向光球半径为
\[
 D_M=R_{K0}\sinh(\chi/R_{K0})
 =\frac c{H_0\sqrt{\Omega_{k0}}}\sinh(\sqrt{\Omega_{k0}}I).
\]
同一光子群今天经过面积$4\pi D_M^2$，每光子能量和到达率各降$1+z$，所以
\[
 J=\frac{PH_0^2\Omega_{k0}}
 {4\pi c^2(1+z)^2\sinh^2[\sqrt{\Omega_{k0}}I]}.
\]
这是从光路、面积和时间能量一起推出，未只引用“与闭情形相同”。

令$\Omega_{k0}\to0$，$\sinh(\sqrt{\Omega_{k0}}I)=\sqrt{\Omega_{k0}}I+O(\Omega_{k0}^{3/2})$，同时趋零项相消，平坦式为
\[
 J=\frac{PH_0^2}{4\pi c^2(1+z)^2I^2}.
\]
对闭宇宙换$\sinh$为$\sin$、$\Omega_k$为绝对值，指定光路且避开焦散。所有版本的单位均为功率／面积；若漏$c^2$，数值单位会错。
\end{solution}

\src{8.286 PS8.2，$\rho_{\rm ms}\propto a^{-3/2}$组分与年龄，(a)--(b)。}
\begin{exercise}
宇宙含非相对论物质、辐射、真空和“神秘物质”，最后一种质量密度随物理体积$V$按$V^{-1/2}$下降。给各组分今天$\Omega_{i0}$及$H_0$。
(a)用$x=a(t)/a(t_0)$写总密度；
(b)写从爆炸至今天的年龄积分。原题指出2011复习库已有官方解，此处独立写出。
\end{exercise}
\begin{solution}
记$\rho_{c0}=3H_0^2/(8\pi G)$、$\Omega_{k0}=1-\Omega_{m0}-\Omega_{r0}-\Omega_{v0}-\Omega_{{\rm ms},0}$，
\[
 \rho(x)=\rho_{c0}
 [\Omega_{m0}x^{-3}+\Omega_{r0}x^{-4}
 +\Omega_{v0}+\Omega_{{\rm ms},0}x^{-3/2}].
\]
曲率不是物质密度的一项，但参与$H^2$。由$\dd t=\dd x/(xH)$，
\[
 t_0=\frac1{H_0}\int_0^1
 \frac{x\,\dd x}
 {\sqrt{\Omega_{r0}+\Omega_{m0}x+\Omega_{k0}x^2+
 \Omega_{{\rm ms},0}x^{5/2}+\Omega_{v0}x^4}}.
\]
这假定当前沿从$x=0$到1单调膨胀的爆炸支，根号全程正；若模型反弹、转向或根本没有爆炸，不能照此选积分下限。连续方程还给该组分$w=-1/2$，它可产生加速，但不是常量真空密度。
\end{solution}

\src{8.286 PS8.3，共同因果过去，实际独算。}
\begin{exercise}
从我们的位置在爆炸瞬间发光，抵达某共动类星体后立即反射，恰在今天收到。题面引Friedman--Kaiser--Gallicchio(2013)的$z=3.65$，同时指定$\Omega_{v0}=0.685,\Omega_{m0}=0.315,\Omega_{r0}=0,H_0=67.3\ \mathrm{km\,s^{-1}\,Mpc^{-1}}$。独立算这个边界，可作数值积分。
\end{exercise}
\begin{solution}
令共形时间$\eta(t)=\int_0^t c\dd t/a$。去程和回程均跨同一共动距离$\chi$，因此在反射事件$t_e$，
\[
 \eta_e=\chi,\quad \eta_0-\eta_e=\chi,\quad
 \chi=\eta_0/2.
\]
接收的红移对应反射时刻$a_e$。平坦模型中$H_0/c$无量纲化后求根条件
\[
 \int_0^{z_e}\frac{\dd z}{\sqrt{.315(1+z)^3+.685}}
 =\frac12\int_0^\infty\frac{\dd z}{\sqrt{.315(1+z)^3+.685}}.
\]
将总积分改为$x=a/a_0$的有限区间$\int_0^1\dd x/\sqrt{.315x+.685x^4}$，避免无限上限误差。实际得总共形积分3.240182741、边界$z_e=3.824781408$。另分别积去程、回程，差小于$1.5\times10^{-13}$。

因此\textbf{题面这组零辐射参数不能精确复现3.65}。实际查\href{https://arxiv.org/html/1305.3943v2}{论文v2式11}，其辐射参数为$9.289\times10^{-5}$，不是零。用课程PS10所给接近值$9.2\times10^{-5}$并保留参数闭合，独算得到$z_e=3.651167$。早时辐射虽今天很小，却改变从爆炸开始的总共形时间，不能在整个积分中省略后仍要求同一数值。论文与课程数值的差别在这里是参数近似的边界，不是因果判据失效。
\end{solution}

\src{8.286 PS8.4，通量表，加分题。}
\begin{exercise}
实际积分平坦$\Omega_{m0}=0.3,\Omega_{v0}=0.7$的$J/(PH_0^2)$，列$z=0.1,0.2,\ldots,1.5$至五位有效数字；也可选开模型$0.3,0.6$，本解选平坦。题面所指“Problem 4”的亮度公式按上下文应为PS8.1。
\end{exercise}
\begin{solution}
采用SI的$c=299792458\ {\rm m\,s^{-1}}$，第一题给$J/(PH_0^2)=1/[4\pi c^2(1+z)^2I^2]$，单位$\mathrm{s^2m^{-2}}$。下表为实际自适应积分，原CSV保留更多数字；若要无量纲表，应再乘$c^2$。
\begin{center}\begin{tabular}{rr|rr|rr}
\toprule
$z$&$J/(PH_0^2)$&$z$&$J/(PH_0^2)$&$z$&$J/(PH_0^2)$\\
\midrule
0.1&$7.6650\,10^{-17}$&0.6&$1.3032\,10^{-18}$&1.1&$2.9384\,10^{-19}$\\
0.2&$1.6908\,10^{-17}$&0.7&$8.9503\,10^{-19}$&1.2&$2.3691\,10^{-19}$\\
0.3&$6.7361\,10^{-18}$&0.8&$6.4501\,10^{-19}$&1.3&$1.9434\,10^{-19}$\\
0.4&$3.4420\,10^{-18}$&0.9&$4.8254\,10^{-19}$&1.4&$1.6180\,10^{-19}$\\
0.5&$2.0236\,10^{-18}$&1.0&$3.7196\,10^{-19}$&1.5&$1.3645\,10^{-19}$\\
\bottomrule\end{tabular}\end{center}
低红移$I\simeq z$，于是通量近$z^{-2}$；较高红移还包含传播几何与两次红移，不能继续拟合单一欧氏幂律。
\end{solution}

\src{8.286 PS8.5，指定历史超新星数据图，加分题。}
\begin{exercise}
选Perlmutter等(1999)或Riess等(2004)公开补充表，画有效星等或距离模数对红移，包含误差棒。叠加至少三个模型：EdS $(\Omega_m,\Omega_v)=(1,0)$，开尘埃$(.3,0)$，平坦含真空$(.3,.7)$（可替代为指定开模型）。各模型星等曲线允许一个整体竖直平移以优化符合。推荐对数红移轴扩至3或线性轴扩至2。解释零点、标准烛光校正及距离模数关系。此解选择Perlmutter公开60行表。
\end{exercise}
\begin{solution}
实际从OCW的\texttt{PSet8\_supp.pdf}逐行解析60颗超新星，保留全部行，不代以8.902 SNLS数据，不声称复现原论文的筛选样本。每颗记录红移、红移误差、有效星等与星等误差。数据是原课程归集的历史测量事实，保留论文署名；不收录原PDF表页或论文图像。

定义$d_L=(H_0/c)D_L$，理论形状$m_{\rm th}(z)=5\log_{10}d_L(z)+b$。未知绝对光度、$H_0$及测光零点都进入同一常数$b$。仅按表中星等误差作示例加权，最优
\[
 b=\frac{\sum_i[m_i-5\log_{10}d_L(z_i)]/\sigma_{m,i}^2}
 {\sum_i1/\sigma_{m,i}^2}.
\]
实际得到$b=24.2422,24.1271,23.9867$，对应三个模型；示例$\chi^2=136.77,112.21,103.18$。这些值没有加入红移误差传播、相关系统误差或原论文筛选，只用于完成题目作图与比较，不是重新发布完整观测似然结论。
\begin{center}\begin{tikzpicture}\begin{axis}[width=.85\textwidth,height=6cm,xmode=log,xlabel={$z$},ylabel={有效星等$m$},xmin=.008,xmax=3,legend pos=north west]
\addplot[only marks,mark size=.9pt,error bars/.cd,y dir=both,y explicit,x dir=both,x explicit] table[x=z,y=magnitude,x error=z_error,y error=magnitude_error,col sep=comma]{qa/early-supernova-derived.csv};\addlegendentry{课程公开60行}
\addplot[black,dashed] table[x=z,y=EdS,col sep=comma]{qa/early-supernova-curves.csv};\addlegendentry{$(1,0)$}
\addplot[gray] table[x=z,y=open,col sep=comma]{qa/early-supernova-curves.csv};\addlegendentry{$(.3,0)$}
\addplot[structurecolor,thick] table[x=z,y=flatLambda,col sep=comma]{qa/early-supernova-curves.csv};\addlegendentry{$(.3,.7)$}
\end{axis}\end{tikzpicture}\end{center}
大星等表示较暗，较远红移处的相对偏差能区分背景膨胀史；单独一个超新星不能解除光度零点与Hubble归一的简并。“有效／校正”星等包含波段与光变宽度的校正，本解沿用表值，没有伪造原始光变拟合。若改选Riess的$\mu$，使用$\mu=5\log_{10}(D_L/{\rm Mpc})+25$，并仍注意表的绝对零点是任意设置。图中比数据更远的曲线是模型外推，没有添加新观测点。
\end{solution}
